%%
%% Test: mb-core — Core Package Loader — v1.1
%% Stage: core refactor (stages 1-4)
%% 
%% Purpose: Verify that all core packages load correctly:
%%   - Mathematics (amsmath, mathtools, amsthm, unicode-math)
%%   - Algorithms (algorithm, algpseudocode)
%%   - Boxes (tcolorbox + minted library)
%%   - Code (minted)
%%   - Bibliography (biblatex)
%%   - Index (imakeidx)
%%   - Layout (geometry, fancyhdr, titlesec, tocloft, setspace, caption)
%%   - Tables (booktabs, array, multirow)
%%   - Hyperref + Cleveref
%%   - pgf-pie (temporary)
%%   - draftwatermark (temporary)
%%   - newcolumntype (temporary)
%% 

\documentclass{matinbook}

\title{تست هسته‌ی MatinBook — نسخه ۱.۱}
\author{تیم MatinBook}
\date{۱۴۰۵}

\begin{document}

\maketitle

\chapter{بررسی هسته‌ی MatinBook}
\label{chap:test-core}

این سند، بارگذاری درست همه‌ی پکیج‌های هسته‌ای را بررسی می‌کند.

%----------------------------------------
\section{ریاضیات}
\label{sec:math}

\subsection{معادلات}

\begin{equation}
    \label{eq:core-test}
    f(x) = \sum_{i=1}^{n} x_i^2 + \int_{0}^{\infty} e^{-t^2} \, dt
\end{equation}

\subsection{نمادها}

مجموعه‌های اعداد: $\mathbb{N} \subset \mathbb{Z} \subset \mathbb{Q} \subset \mathbb{R} \subset \mathbb{C}$

\subsection{محیط‌های ریاضی}

\begin{align}
    a &= b + c \\
    d &= e + f
\end{align}

\subsection{قضیه}

\begin{theorem}[قضیه‌ی نمونه]
    برای هر عدد حقیقی $x$، داریم $x^2 \geq 0$.
\end{theorem}

\begin{proof}
    چون مربع هر عدد حقیقی نامنفی است، نتیجه حاصل می‌شود.
\end{proof}

%----------------------------------------
\section{الگوریتم‌ها}
\label{sec:algorithms}

\begin{latin}
\begin{algorithm}
\caption{محاسبه‌ی فاکتوریل}
\begin{algorithmic}[1]
    \Require عدد صحیح $n \geq 0$
    \Ensure $n!$
    \State $result \gets 1$
    \For{$i \gets 2$ \textbf{to} $n$}
        \State $result \gets result \times i$
    \EndFor
    \State \Return $result$
\end{algorithmic}
\end{algorithm}
\end{latin}

%----------------------------------------
\section{جعبه‌ها}
\label{sec:boxes}

\begin{tcolorbox}[title=جعبه‌ی نمونه]
این یک جعبه‌ی ساده است که با \texttt{tcolorbox} ساخته شده.
\end{tcolorbox}

%----------------------------------------
\section{کد}
\label{sec:code}

\begin{latin}
\begin{minted}{python}
def factorial(n):
    """Compute n! recursively."""
    if n <= 1:
        return 1
    return n * factorial(n - 1)
\end{minted}
\end{latin}

%----------------------------------------
\section{جدول}
\label{sec:tables}

\begin{table}[h]
\centering
\caption{جدول نمونه}
\label{tab:core-test}
\begin{tabular}{L{3cm} C{2cm} R{3cm}}
\toprule
\textbf{ستون ۱} & \textbf{ستون ۲} & \textbf{ستون ۳} \\
\midrule
مقدار ۱ & ۱۰ & مقدار ۳ \\
مقدار ۲ & ۲۰ & مقدار ۴ \\
\bottomrule
\end{tabular}
\end{table}

%----------------------------------------
\section{ارجاعات}
\label{sec:references}

ارجاع به معادله: \cref{eq:core-test}

ارجاع به جدول: \cref{tab:core-test}

ارجاع به فصل: \cref{chap:test-core}

%----------------------------------------
\section{نتیجه‌گیری}
\label{sec:conclusion}

اگر این سند بدون خطا کامپایل شود:

\begin{itemize}
    \item ✅ ریاضیات (\texttt{amsmath}, \texttt{mathtools}, \texttt{unicode-math}) کار می‌کند
    \item ✅ قضیه (\texttt{amsthm}) کار می‌کند
    \item ✅ الگوریتم (\texttt{algorithm}, \texttt{algpseudocode}) کار می‌کند
    \item ✅ جعبه‌ها (\texttt{tcolorbox}) کار می‌کنند
    \item ✅ کد (\texttt{minted}) کار می‌کند
    \item ✅ جدول (\texttt{booktabs}, \texttt{array}, \texttt{multirow}) کار می‌کند
    \item ✅ ارجاعات (\texttt{hyperref}, \texttt{cleveref}) کار می‌کنند
    \item ✅ \texttt{newcolumntype}های سفارشی کار می‌کنند
\end{itemize}

\textbf{وضعیت تست:} قبول

\end{document}
